\section{The Angelic Substance}\label{sec:substance}

After treating the distinction of creatures in general, Aquinas turns to
the distinction of corporeal and spiritual creature (\latin{de
distinctione corporalis et spiritualis creaturae}) and takes first the
purely spiritual creature, which in Sacred Scripture is named angel
(\latin{de creatura pure spirituali, quae in Scriptura sacra Angelus
nominatur}), then the purely corporeal creature, and then man, composed of
both (\ST{I q.~50 pr.}). The angels themselves are considered under four
heads: what pertains to their substance, to their intellect, to their will,
and to their creation (\latin{primo de his quae pertinent ad eorum
substantiam; secundo~\dots\ ad eorum intellectum; tertio~\dots\ ad eorum
voluntatem; quarto~\dots\ ad eorum creationem}). The substance is taken
first absolutely (\latin{absolute}, q.~50) and then in comparison with
corporeal things (\latin{per comparationem ad corporalia}): bodies, place, and local
motion (qq.~51--53, \S\ref{sec:bodies}). Intellect follows in qq.~54--58
(\S\ref{sec:knowledge}), will and love in qq.~59--60 (\S\ref{sec:will}),
and the angels' coming-to-be in qq.~61--64, which the prologue of q.~61
divides into their production in natural being, their perfection in grace
and glory, and the evil of some (\S\S\ref{sec:creation},~\ref{sec:fall}).
The angels' part in the divine government (qq.~106--114) belongs to a later
treatise of the \emph{Prima pars}.

Every later question rests on q.~50. The angel is a created intellect
without a body and without matter, a form subsisting in its own act of
being and yet not identical with that being; from this follow the vast
number of the angels, the distinction of each from each by species, and
their natural incorruptibility.

\subsection{The Substance of the Angels (\ST{I q.~50})}\label{sec:q50}

The five articles ask whether there is any creature wholly spiritual and
incorporeal; whether such a creature is composed of matter and form; how
many angels there are; how they differ from one another; and whether they
are incorruptible (\ST{I q.~50 pr.}).

\paragraph{Wholly incorporeal creatures (a.~1).} The \emph{sed contra} is
Ps 103[104]:4, ``Who makest thy angels spirits'' (Douay). Aquinas argues
from the end of creation. God intends in creatures their good, which lies
in likeness to himself, and the likeness of an effect is perfect when it
imitates the cause in that by which the cause acts; God acts by intellect
and will, so ``the perfection of the universe requires that there should be
intellectual creatures.'' But ``intelligence cannot be the action of a
body, nor of any corporeal faculty; for every body is limited to `here' and
`now'\,'' (\latin{omne corpus determinatur ad hic et nunc}). The universe is
therefore incomplete without an incorporeal creature. The contrary error
came from not distinguishing sense from intellect: the ancients admitted
only what imagination grasps, and ``thence came the error of the
Sadducees, who said there was no spirit'' (Acts 23:8). The first objection
quotes Damascene that the angel is corporeal and material compared with
God. Damascene writes: ``For all that is compared
with God Who alone is incomparable, we find to be dense and material. For
in reality only the Deity is immaterial and incorporeal'' (\emph{De fide
orthodoxa} II.3). Aquinas answers that the angels stand between God and
bodies, and ``the medium compared to one extreme appears to be the other
extreme,'' as tepid seems cold beside heat (ad~1). Damascene's ``ever
mobile'' (\latin{semper mobilis}) refers to the unceasing act of intellect
and will (ad~2). Ambrose's principle that every creature is bounded answers
itself: ``To be circumscribed by local limits belongs to bodies only;
whereas to be circumscribed by essential limits belongs to all creatures,
both corporeal and spiritual'' (ad~3). Ambrose's own sentence argues the
same distinction in the service of the Spirit's divinity: ``every creature
is confined within certain limits of its own nature,'' and invisible
operations that ``cannot be circumscribed by place and bounds, yet are
closed in by the property of their own substance'' (\emph{De Spiritu
Sancto} I.7.81).

\paragraph{No matter in the angels (a.~2).} The \emph{sed contra} is
Dionysius, \emph{Divine Names} 4, that the first creatures are understood
to be immaterial as they are incorporeal; in Parker's rendering the
intelligible essences ``are conceived of as incorporeal and immaterial''
(\emph{DN}~4.1). Aquinas names the adversary: Avicebron, in the \emph{Fons
vitae}, held one universal matter of spiritual and corporeal things,
because he assumed that whatever the intellect distinguishes is really
distinct. Aquinas answers in two steps. First, ``one glance is enough to
show that there cannot be one matter of spiritual and of corporeal
things'': a single matter could receive the two forms only in different
parts, and matter has parts only under quantity, so spiritual matter would
be quantified. Second, no intellectual substance has any matter at all,
because operation follows the mode of substance, and ``to understand is an
altogether immaterial operation,'' as its object shows: forms are
intelligible in so far as they are abstracted from matter. Avicebron's root
mistake is answered by the principle that ``the intellect does not
apprehend things according to their mode, but according to its own mode''
(\latin{intellectus non apprehendit res secundum modum rerum, sed secundum
modum suum}). The replies carry the metaphysics of the whole treatise. In
immaterial things genus and difference are taken from one and the same
reality under two conceptions (ad~1); the intellect receives forms as
forms, not as matter receives them (ad~2). Above all: ``Although there is
no composition of matter and form in an angel, yet there is act and
potentiality'' (ad~3). The angel is a \latin{forma subsistens} related to
its own \latin{esse} as potency to act, so that it is composed, in
Boethius's terms, of \latin{quod est} and \latin{esse}: ``For `what is,' is
the form itself subsisting; and the existence itself is whereby the
substance is; as the running is whereby the runner runs.'' In God alone
\latin{esse} and \latin{quod est} are the same; ``hence God alone is pure
act.'' The angelic form is finite in being and ``infinite'' only in not
being received into matter; hence the \emph{Liber de causis}: ``intelligence
is finite from above'' and ``infinite from below'' (ad~4).

Bonaventure determines the question the other way. The angel is mutable,
passible, individuated and limited, and essentially composite in its own
nature; therefore \latin{non video causam nec rationem, quomodo defendi
potest, quin substantia Angeli sit composita ex diversis naturis}, and
\latin{illa positio videtur verior esse, scilicet quod in Angelo sit
compositio ex materia et forma} (\emph{In II Sent.} d.~3 p.~1 a.~1 q.~1).
A creature, he replies, is not pure act and so must have possibility;
being mutable, it must have a foundation; being limited and in a genus, it
must have composition. The authorities who deny matter to spirits speak of
matter \latin{appropriate}, as the principle of undergoing change (ad~1,
ad~3). Aquinas concedes the name in the wide sense and refuses it as
improper: if every potency is called matter and every act form, there is
matter and form in spiritual substance, \latin{sed tamen hoc non est
proprie dictum secundum communem usum nominum} (\emph{De spiritualibus
creaturis} a.~1). Both masters hold that the angel is composed of potency
and act. They divide over whether that potency is matter, and the
dispute runs on in the schools (\S\ref{sec:dq-matter}).

\paragraph{The multitude of the angels (a.~3).} Dan 7:10 is the \emph{sed
contra}: ``thousands of thousands ministered to him, and ten thousand times
a hundred thousand stood before him'' (Douay). Aquinas sets aside Plato,
who counted the separate substances by the species of sensible things;
Aristotle, who counted them by the first motions of the heavens; and
``Rabbi Moses'' (Maimonides), who kept Aristotle's number for the
immaterial substances and read Scripture's other angels as prophets or
natural powers. Against the last he replies that ``it is, however, quite
foreign to the custom of the Scriptures for the powers of irrational things
to be designated as angels.'' His determination is that ``the angels, even
inasmuch as they are immaterial substances, exist in exceeding great
number, far beyond all material multitude.'' He cites \emph{CH}~14, which
in Parker's translation reads: ``many are the blessed hosts of the
supermundane minds, surpassing the weak and contracted measurement of our
material number.'' The reason is again the perfection of the universe: the
more perfect things are, the more abundantly God creates them. In bodies
that abundance shows as magnitude, the heavens dwarfing the sphere of
generation; in incorporeal things it shows as multitude. The replies secure
the kind of number meant. Angelic number is not discrete quantity but the
transcendental multitude that follows distinction of forms (ad~1). The
nature nearest to God has least multitude in its composition, not in its
subjects (ad~2). The separate substances do not exist for the sake of the
heavens, so Aristotle's count is only probable (ad~3). Their multiplication
follows ``the divine wisdom devising the various orders of immaterial
substances'' (ad~4).

\paragraph{Each angel a species (a.~4).} The \emph{sed contra} draws on
the Dionysian hierarchy. Within one species there is no ``first'' and
``second,'' yet even within one order there are first, middle, and last
angels. Dionysius has ``every Hierarchy distributed into first, and middle,
and last Powers'' (\emph{CH}~10.2). Aquinas records three opinions: that
all spiritual substances, souls included, are of one species; that all
angels are of one species, though souls are not; and that the angels of
one hierarchy or one order are of one species. He rejects all three. Things
one in species and many in number agree in form and differ by matter; since
the angels have no matter, ``it is impossible for two angels to be of one
species,'' just as a separated whiteness or humanity could only be one.
Even if they had matter, it would be distinguished by diversity of
potencies rather than by quantity, and that would make them differ in
genus. The replies explain the ranking this produces: ``all the angels
differ in species according to the diverse degrees of intellectual
nature,'' as fire differs from air, by forms of different grade, not by
intensity of one form (ad~1--2). ``The good of the species preponderates
over the good of the individual,'' so multiplying species is the better
way to multiply angels (ad~3). The agent intends specific rather than
numerical multitude (ad~4). Damascene leaves the question to God: ``Whether
they are equals in essence or differ from one another we know not. God,
their Creator, Who knoweth all things, alone knoweth'' (\emph{De fide
orthodoxa} II.3). The dissent of the
Franciscan school is recorded early. A late thirteenth-century index of
places where Bonaventure and Thomas disagree on the second book of the
\emph{Sentences}, printed by the Quaracchi editors, lists \latin{Utrum
omnes Angeli differant specie} with \latin{Thomas, quod sic; Bonaventura,
quod non omnes} (Bonaventure, \emph{Opera omnia} II, Quaracchi 1885,
preface, index no.~11).

\paragraph{Incorruptible by nature (a.~5).} The \emph{sed contra} is again
\emph{DN}~4. In Parker's version the intelligible essences ``have their
life, continuous and undiminished, purified from all corruption and death
and matter, and generation'' (\emph{DN}~4.1). Aquinas's reason: nothing
corrupts except by the separation of form from matter, and to be belongs to
a form by itself, as roundness belongs to a circle. A composite ceases to
be when its form leaves its matter, but ``if the form subsists in its own
being, as happens in the angels~\dots\ it cannot lose its being.''
Therefore ``the angel's immateriality is the cause why it is incorruptible
by its own nature'' (\latin{ipsa igitur immaterialitas Angeli est ratio
quare Angelus est incorruptibilis secundum suam naturam}). Intellectual
operation gives the sign of it, since its object is above time. The first
objection is Damascene's own sentence: ``It is immortal, not by nature but
by grace'' (\emph{De fide orthodoxa} II.3). Aquinas reads it of ``perfect
immortality, which includes complete immutability,'' which the angels have
only by grace, as q.~62 will show (ad~1; \S\ref{sec:q62}). The third
objection cites Gregory, who writes: ``For all things
were made out of nothing, and their being would again go on into nothing,
except the Author of all things held it by the hand of governance''
(\emph{Moralia} XVI.37.45). Aquinas grants the dependence and denies that
it is corruptibility: a necessary being may have a cause of its necessity,
and Gregory's text shows only ``that the nature of the angels is dependent
upon God as its cause.'' A thing is corruptible when it has within it a
principle of corruption, a contrariety, or ``at least the potentiality of
matter'' (ad~3). Plato's gods of the \emph{Timaeus}, dissoluble by nature
and preserved by the maker's will, are the heavenly bodies (ad~2).

\angeldossier{\ST{I q.~50 aa.~1--5}; parallels \emph{Summa contra
gentiles} II.46, 49--55; \emph{De spiritualibus creaturis} aa.~1, 8;
\emph{De ente et essentia} c.~4; \emph{De substantiis separatis} cc.~5--8.}
 {Defined: God created from nothing at the beginning of time a spiritual
 creature distinct from the corporeal, the angelic (Lateran IV,
 \emph{Firmiter}, DS~800: \latin{utramque de nihilo condidit creaturam,
 spiritualem et corporalem, angelicam videlicet et mundanam}; Vatican~I,
 \emph{Dei Filius} ch.~1). Church teaching: the angels are spiritual,
 non-corporeal beings, whose existence is ``a truth of faith'' (CCC~328),
 ``personal and immortal creatures'' (CCC~330) (aa.~1, 5). Common
 teaching: the angels exist in a number beyond reckoning (a.~3; Dan 7:10;
 \emph{CH}~14). Thomist position: no matter in the angels, only the
 composition of \latin{esse} and \latin{quod est} (a.~2); each angel its
 own species (a.~4); incorruptibility belongs to the angelic nature as
 such (a.~5).}
 {Ps 103[104]:4, Ps 148:2,~5, Dan 7:10, Acts 23:8; Dionysius, \emph{DN}~4.1 (\emph{sed contra} of aa.~2 and~5, obj.~4 of a.~3), \emph{CH}~10.2
 and~14; Damascene, \emph{De fide orthodoxa} II.3 (objections of aa.~1
 and~5); Ambrose, \emph{De Spiritu Sancto} I.7.81; Gregory, \emph{Moralia}
 XVI.37.45; Boethius, \emph{De Trinitate} (obj.~2 of a.~2); \emph{Liber de causis}; Aristotle; Maimonides and Plato as
 reported by Aquinas; Lateran IV (DS~800); Vatican~I, \emph{Dei Filius};
 CCC~328--330.}
 {Avicebron, \emph{Fons vitae}, one universal matter of spiritual and
 corporeal things (a.~2; \emph{De sub.\ sep.} cc.~5--8); Bonaventure,
 \emph{In II Sent.} d.~3 p.~1 a.~1 q.~1, matter taken broadly in the angel
 (a.~2); the three unnamed opinions on angels sharing a species, and the
 Franciscan dissent noted in the thirteenth-century index (a.~4);
 Damascene, uncertain whether angels differ in essence (a.~4), and
 immortal ``not by nature but by grace'' (a.~5); Aristotle and Maimonides
 on the number of separate substances (a.~3).}
